\chapter{A Concrete Naming Certificate for $\RegularCauchyModulusDominationUp$}
\label{ch:concrete-instances-regular-cauchy-modulus-domination-namecert}
\origin{human}

Following Bishop and Bridges, the $\RegularCauchyModulusDominationUp$ packet records the finite domination route by which one regular-Cauchy modulus may be replaced by a cofinal dominating modulus without changing the displayed tail-window readback. It sits after the completion-modulus packet of \autoref{ch:concrete-instances-bishopcompletionmodulus-namecert}, the Cauchy threshold surface of \autoref{ch:concrete-instances-cauchymodulus-namecert}, the finite observation windows of \autoref{ch:concrete-instances-streamname-namecert}, the dyadic tolerance ledger of \autoref{ch:concrete-instances-dyadicratcore-namecert}, the regular rational handoff of \autoref{ch:concrete-instances-regseqrat-namecert}, and the terminal real-name boundary of \autoref{ch:concrete-instances-real-namecert}. The packet names a finite L10 route for changing modulus representatives along a displayed tail comparison; it does not select limits, quotient regular Cauchy names, countable-choice schedules, host real equality, or ambient completeness.

\begin{definition}[Regular Cauchy modulus domination carrier]
\label{def:regular-cauchy-modulus-domination-carrier}
A $\RegularCauchyModulusDominationUp$ carrier is a finite $\BHist$ packet
$$
K=(M_0,M_1,D,W_0,W_1,T,R,E,H,C,P,N)
$$
with the following displayed rows:
\begin{enumerate}
\item $M_0$ is the source $\CauchyModulusUp$ row attached to the original regular-Cauchy tail request.
\item $M_1$ is the dominating $\CauchyModulusUp$ row whose displayed threshold is used by the replacement route.
\item $D$ is the $\DyadicRatCoreUp$ domination ledger, certifying that $M_1$ dominates $M_0$ on the displayed tail and tolerance scale.
\item $W_0$ and $W_1$ are the corresponding $\StreamNameUp$ finite windows read from the source modulus and the dominating modulus.
\item $T$ is the cofinal tail comparison row, recording that every source-tail read used by $M_0$ is met by a dominating-tail read through $M_1$.
\item $R$ is the $\RegSeqRatUp$ handoff row obtained from the dominated window comparison and the dyadic ledger.
\item $E$ is the $\RealUp$ seal row reached only after the two moduli, the domination ledger, both windows, the cofinal tail row, and the regular readback have been displayed.
\item $H$ stores componentwise $\hsame$ transport for the two modulus rows, domination ledger, windows, cofinal tail comparison, readback, seal, replay, provenance, and naming rows.
\item $C$ records displayed $\Cont$ replay of the route
$$
M_0,M_1,D \longmapsto W_0,W_1,T \longmapsto R \longmapsto E .
$$
\item $P$ is the $\Pkg$ provenance over $\RegularCauchyModulusDominationUp$, $\BishopCompletionModulusUp$, $\CauchyModulusUp$, $\StreamNameUp$, $\DyadicRatCoreUp$, $\RegSeqRatUp$, $\RealUp$, $\BHist$, $\hsame$, $\Cont$, and $\NameCert$ rows.
\item $N$ is the local $\NameCert_{\RegularCauchyModulusDominationUp}$ row.
\end{enumerate}
The classifier compares accepted carriers only through $M_0,M_1,D,W_0,W_1,T,R,E,H,C,P,N$. It has no coordinate for an infinite modulus choice function, selected limit, quotient stream equality, countable choice, host real equality, or ambient $\RealUp$ completeness theorem.
\end{definition}
\leanstmt{BEDC.Derived.RegularCauchyModulusDominationUp}

\begin{theorem}[Regular Cauchy modulus domination NameCert obligations]
\label{thm:regular-cauchy-modulus-domination-namecert-obligations}
For every accepted $\RegularCauchyModulusDominationUp$ carrier $K=(M_0,M_1,D,W_0,W_1,T,R,E,H,C,P,N)$, the local $\NameCert_{\RegularCauchyModulusDominationUp}$ surface consists exactly of carrier admission for the displayed packet; source and dominating modulus exposure; dyadic domination through $D$; paired StreamName window exposure through $W_0,W_1$; cofinal tail comparison through $T$; RegSeqRat handoff through $R$; RealUp seal non-escape through $E$; componentwise $\hsame$ transport through $H$; displayed $\Cont$ replay through $C$; $\Pkg$ provenance through $P$; and local naming through $N$.
\end{theorem}
\begin{proof}
Project the carrier. The tuple in \autoref{def:regular-cauchy-modulus-domination-carrier} lists exactly the two modulus rows, the dyadic domination ledger, the two window rows, the cofinal tail comparison, the regular readback, the Real seal, and the structural rows.

The domination obligation is forced by the order of the rows. A consumer first reads $M_0$ and $M_1$, then reads $D$ as the finite certificate that the displayed $M_1$ threshold dominates the displayed $M_0$ threshold on the inspected tail.

The handoff rows are downstream. The windows $W_0,W_1$ and the cofinal comparison $T$ supply the only route to $R$, and $E$ is available only after $R$. Transport, replay, provenance, and naming preserve the same packet, so no unlisted modulus selector or completion principle can be exported.
\end{proof}

\begin{theorem}[Regular Cauchy modulus domination tail stability]
\label{thm:regular-cauchy-modulus-domination-tail-stability}
Let $K=(M_0,M_1,D,W_0,W_1,T,R,E,H,C,P,N)$ be an accepted $\RegularCauchyModulusDominationUp$ carrier. Any later tail request transported by the displayed $\hsame$ rows and still certified by the domination ledger $D$ reads the same cofinal comparison row $T$, the same dominating window route $W_1$, and the same $\RegSeqRatUp$ handoff row $R$. Thus replacing $M_0$ by the dominating modulus $M_1$ is stable for every Real-facing read that remains inside the displayed tail.
\end{theorem}
\begin{proof}
Start with the transported tail request. Componentwise $\hsame$ transport preserves the visible modulus, ledger, window, cofinal comparison, readback, and seal rows rather than creating a second tail source.

Read the dyadic domination ledger next. Since $D$ is still the certificate for the transported request, the route from the source window to the dominating window must pass through $T$ and then through $W_1$.

Finally replay the continuation route. The row $R$ is downstream from the two windows and $T$, and the Real seal $E$ is downstream from $R$. Hence every transported Real-facing read stays on the original dominated tail route.
\end{proof}

\begin{theorem}[Regular Cauchy modulus domination Real seal route]
\label{thm:regular-cauchy-modulus-domination-real-seal-route}
Every accepted $\RegularCauchyModulusDominationUp$ carrier routes a Real-completion consumer through the finite chain
$$
M_0,M_1,D \longmapsto W_0,W_1,T \longmapsto R \longmapsto E .
$$
The terminal $\RealUp$ seal records that the dominated cofinal tail window has already been repacked as a $\RegSeqRatUp$ handoff. It exports no selected real limit, quotient regular Cauchy equality, hidden modulus choice function, countable choice, host real equality, or ambient completeness principle.
\end{theorem}
\begin{proof}
The displayed $\Cont$ row supplies the chain. A Real-completion consumer must first expose the two modulus rows and the domination ledger, then the paired windows and cofinal tail comparison.

The regular rational handoff is the next visible row. By the carrier definition, $R$ is obtained only after $W_0,W_1,T$ have been displayed, so the handoff cannot be read from an unregistered modulus or a hidden tail selector.

The Real seal is terminal. Since $E$ appears after $R$ and the carrier contains no selected limit, quotient equality, choice schedule, host equality, or completeness coordinate, the seal exports only the finite dominated-tail route.
\end{proof}

\falsifiablePrediction{An accepted $\RegularCauchyModulusDominationUp$ consumer that reaches a $\RealUp$ seal without displaying $M_0$, $M_1$, the dyadic domination ledger $D$, both $\StreamNameUp$ windows, the cofinal tail comparison $T$, the $\RegSeqRatUp$ handoff, componentwise $\hsame$ transport, displayed $\Cont$ replay, $\Pkg$ provenance, and local $\NameCert$ rows refutes the packet.}

\independenceWitness{bishopcompletionmodulus, cauchymodulus, streamname, dyadicratcore, regseqrat, real: no listed sibling is bijective with $\RegularCauchyModulusDominationUp$ because this packet jointly carries two modulus rows, a domination ledger, paired windows, cofinal tail comparison, regular readback, Real sealing, transport, replay, provenance, and local naming.}

\closureat{RegularCauchyModulusDominationUp}{seedStr}

\begin{closurestatus}{\RegularCauchyModulusDominationUp}
  \constructivestory{A $\RegularCauchyModulusDominationUp$ packet is generated from two CauchyModulus rows, a DyadicRatCore domination ledger, paired StreamName windows, a cofinal tail comparison row, a RegSeqRat handoff, a RealUp seal, componentwise hsame transport, displayed Cont replay, Pkg provenance, and a local NameCert row.}
  \theoryclosure{\seedClosure}
  \scopeclosed{\autoref{def:regular-cauchy-modulus-domination-carrier}, \autoref{thm:regular-cauchy-modulus-domination-namecert-obligations}, \autoref{thm:regular-cauchy-modulus-domination-tail-stability}, and \autoref{thm:regular-cauchy-modulus-domination-real-seal-route} seed the finite domination route between source and dominating regular-Cauchy moduli over displayed tail windows.}
  \formalstatus{\unformalizedV}
  \bridgestatus{none}
  \notclaimed{This chapter does not close selected limits, quotient regular Cauchy equality, infinite modulus choice, countable choice, host real equality, ambient Real completeness, Lean verification, or bridge checking.}
  \upgradepath{Obligation closure requires checked carrier admission, dyadic domination, paired window exposure, cofinal tail comparison, RegSeqRat handoff, Real seal non-escape, transport, replay, provenance, and local naming for BEDC.Derived.RegularCauchyModulusDominationUp.}
\end{closurestatus}
